%This is a template file for use of iopjournal.cls

\documentclass{iopjournal}

% Packages
\usepackage{amsfonts}
\usepackage{amsmath}

% Options
% 	[anonymous]	Provides output without author names, affiliations or acknowledgments to facilitate double-anonymous peer-review

\begin{document}

\articletype{Paper} %	 e.g. Paper, Letter, Topical Review...

\title{Conformalized Quantum DeepONet Ensembles for Scalable Operator Learning with Distribution-Free Uncertainty}

\author{Purav Matlia$^{1, *}$\orcid{0009-0006-3944-1665}, Christian Moya$^2$\orcid{0000-0003-0180-9285} and Guang Lin$^{2,3}$\orcid{0000-0002-0976-1987}}

\affil{$^1$Department of Computer Science, Purdue University, West Lafayette, IN 47907, United States of America}

\affil{$^2$Department of Mathematics, Purdue University, West Lafayette, IN 47907, United States of America}

\affil{$^3$School of Mechanical Engineering, Purdue University, West Lafayette, IN 47907, United States of America}

\affil{$^*$Author to whom any correspondence should be addressed.}

\email{pmatlia@purdue.edu, cmoyacal@purdue.edu, and guanglin@purdue.edu}

\keywords{Quantum Machine Learning, Operator Learning, Uncertainty Quantification, Conformal Prediction, Dynamical Systems}

\begin{abstract}
Operator learning enables fast surrogate modeling of high-dimensional dynamical systems, but existing approaches face two fundamental limitations: quadratic inference complexity and unreliable uncertainty quantification in safety-critical settings.
We propose Conformalized Quantum DeepONet Ensembles, a framework that addresses both challenges simultaneously. By leveraging Quantum Orthogonal Neural Networks (QOrthoNNs), we reduce operator inference complexity from $\mathcal{O}(n^2)$ to $\tilde{\mathcal{O}}(n)$, enabling scalable evaluation over fine discretizations. To provide rigorous uncertainty quantification, we combine ensemble-based epistemic modeling with adaptive conformal prediction, yielding distribution-free coverage guarantees. A key challenge in ensembling is that naïve parallelism scales hardware resources linearly with the number of models. We resolve this by using Superposed Parameterized Quantum Circuits (SPQCs), which compress multiple ensemble members into a single circuit and enable simultaneous multi-model execution. Experiments on synthetic partial differential equations and real-world power system dynamics demonstrate that our approach achieves accurate predictions while maintaining calibrated uncertainty under realistic quantum noise. These results establish a practical pathway toward scalable, uncertainty-aware operator learning in quantum machine learning.
\end{abstract}




\section{Introduction}
Evaluating dynamical systems across large sets of hypothetical scenarios is central to scientific computing and safety-critical applications such as power systems. While neural operator methods such as DeepONet provide fast surrogate models for these tasks, they face two fundamental challenges: (i) scalability, due to the quadratic $\mathcal{O}(n^2)$ cost of dense neural layers, and (ii) reliability, as standard uncertainty quantification methods either lack rigorous guarantees or incur prohibitive computational overhead. These limitations are particularly acute in settings that require both fine discretization and trustworthy predictions.

We address these challenges by combining quantum operator learning with distribution-free uncertainty quantification. Our key observation is that while quantum architectures can reduce inference complexity, uncertainty estimation via ensembles introduces a competing resource bottleneck, as naïve parallelism scales qubit requirements linearly. Resolving this tension requires a joint design of scalable inference, efficient ensembling, and calibrated uncertainty.

We make the following contributions:
\begin{itemize}
\item \textbf{Scalable quantum operator inference}. We build on Quantum DeepONet architectures based on Quantum Orthogonal Neural Networks to achieve scalable operator learning with reduced inference complexity from $\mathcal{O}(n^2)$ to $\tilde{\mathcal{O}}(n)$.
\item \textbf{Resource-efficient hybrid and superposed ensembling}. We mitigate the hardware overhead of ensembling through two architectural strategies. First, we introduce a hybrid classical-quantum architecture that reduces qubit requirements and reduces the effect of quantum hardware noise. Second, we deploy Superposed Parameterized Quantum Circuits (SPQCs) to encode all ensemble predictions within a single circuit. Unlike sequential execution—which requires multiple distinct state preparations, data encodings, and measurements—SPQCs execute these steps exactly once. The unique architecture of the Quantum DeepONet allows these strategies to complement each other.
\item \textbf{Distribution-free uncertainty quantification}. We integrate adaptive conformal prediction with ensemble outputs to provide rigorous, finite-sample coverage guarantees.
\item \textbf{Robust evaluation under realistic noise}. We demonstrate accurate and calibrated performance across synthetic PDEs and real-world power system tasks under hardware-calibrated quantum noise models.
\end{itemize}

Unlike prior work that studies quantum operator learning or conformal prediction in isolation, our framework addresses their interaction under resource constraints. In particular, while ensembles are essential for epistemic uncertainty, their naïve parallel implementation introduces a prohibitive qubit overhead. We address this tension by jointly designing the model architecture, ensemble execution, and calibration procedure, enabling uncertainty-aware operator learning without sacrificing computational efficiency.

Traditional numerical simulations accurately analyze dynamical systems governed by high-dimensional differential equations \cite{dahlquist2008numerical, butcher2016numerical, liu2019mcsimulation}. However, repeatedly evaluating these simulators for thousands of hypothetical perturbations is computationally expensive. This cost motivates using learned surrogate operators \cite{lu2021deeponet, li2023neuraloperators, azizzadenesheli2024neuraloperators, wen2022neuraloperator} to approximate the solution map between infinite-dimensional function spaces directly. While these neural operator architectures provide flexibility, they introduce new challenges regarding scalability and uncertainty quantification \cite{psaros2023uq, zou2024operatoruq, azizzadenesheli2024neuraloperators, xiao2025qdon, kerenidis2022qorthonn, landman2022medicalqorthonn}.

Deep Operator Networks (DeepONet) \cite{lu2021deeponet, lu2022deeponetcomparison, lu2022multifidelitydeeponet, lin2023responsedeeponet} approximate solution maps by decomposing them into branch and trunk sub-networks. However, classical hidden layers within these networks rely on dense matrix-vector multiplications that incur an $O(n^2)$ computational cost \cite{kerenidis2022qorthonn, landman2022medicalqorthonn, goodfellow2016deeplearning}. This quadratic scaling is prohibitive for Monte Carlo uncertainty quantification \cite{zhang2021mccost, milanovic2017mccost, ye2023mccost}, which requires evaluating $10^4$ timesteps to capture fine-scale transient dynamics. Overcoming this bottleneck requires replacing classical linear layers with a sub-quadratic inference mechanism.

Recent quantum advancements introduce scalable architectures like the Quantum Fourier Neural Operator (QFNO) \cite{jain2024qfno} and Quantum DeepONet \cite{xiao2025qdon}. By replacing $\mathcal{O}(n^2)$ classical matrix multiplications with parameterized Reconfigurable Beam Splitter (RBS) circuits operating in the unary subspace \cite{kerenidis2022qorthonn, landman2022medicalqorthonn, xiao2025qdon}, Quantum DeepONets reduce forward-pass complexity to $\tilde{O}(n)$. This efficiency enables the fine discretization necessary to capture high-frequency transients. However, because quantum measurements produce probabilistic outputs requiring multiple shots, a single Quantum DeepONet yields only a noisy point prediction. Consequently, these accelerated predictions alone fail to provide the epistemic uncertainty estimates required for safety-critical deployments.

To capture true epistemic uncertainty, we rely on model diversity in ensembles and conformal prediction. Although Bayesian Neural Networks \cite{neal1996bnn, goan2020bnn} or Monte Carlo dropout \cite{gal2016dropout} traditionally model this uncertainty, they lack the distribution-free coverage guarantees required for safety-critical tasks. Conformalized deep ensembles \cite{angelopoulous2023conformal, moya2025conformal} resolve this by providing rigorous mathematical guarantees and model-agnostic epistemic uncertainty. While methods like Quantum Conformal Prediction (QCP) \cite{park2023qcp} attempt to bypass ensemble costs by leveraging the inherent measurement shots of a single quantum circuit as a probabilistic predictor, this approach fails for quantum DeepONets as shot variance captures only hardware-level noise rather than true structural epistemic diversity. However, naïvely replicating $L$ parallel quantum DeepONets scales qubit requirements with $L$. This linear scaling reintroduces a hardware bottleneck that undermines the scalability required for massive Monte Carlo evaluations. Consequently, the central challenge is preserving ensemble-based epistemic diversity without proportionally increasing resource costs.

We address these bottlenecks by introducing a unified quantum ensemble framework. This approach integrates scalable quantum operator inference, model-level ensembling, and resource-efficient execution, via hybrid classical–quantum architectures and Superposed Parameterized Quantum Circuits \cite{patapovich2025spqc}, alongside conformal calibration to guarantee distribution-free coverage. 

The rest of the paper is organized as follows. Following a review of previous work in Section \ref{sec:previous_work}, we formulate the problem in Section \ref{sec:problemformulation}, Section \ref{sec:method} details the quantum circuit design and ensembling methodology, Section \ref{sec:uq} establishes the conformal prediction framework, while Section \ref{sec:experiments} evaluates the framework's predictive accuracy and uncertainty calibration. Finally, we discuss limitations in Section \ref{sec:limitations}. For clarity, a summary of nomenclature is provided in Appendix \ref{sec:nomenclature}, which readers may consult throughout.




\begin{figure}[t]

\begin{center}
\includegraphics[width=\columnwidth]{figures/main_figure.pdf}
\end{center}

\caption{Overview of the conformalized quantum DeepONet ensemble framework. (A) quantum DeepONet ensemble: an ensemble of operators decomposed into two subnetworks implemented using quantum orthogonal neural networks. (B) ensemble output: predictions are aggregated to compute mean prediction $\mu(t)$ and standard deviation $\sigma(t)$ for uncertainty. (C) conformal safety: using calibration data $\mathcal{D}_{cal}$, nonconformity scores define threshold $\hat{q}$ producing pointwise prediction intervals $C_{\alpha}(\mathbf{u}, t)$ with coverage $1-\alpha$.}

\label{fig:main_framework}
\end{figure}



\section{Previous Work}
\label{sec:previous_work}

\paragraph{Operator learning and scientific machine learning.}
The scientific machine learning community has introduced various data-driven frameworks for approximating the dynamic response of systems governed by ordinary and partial differential equations \cite{karniadakis2021physicsinformedml, baker2019scimlreport}. Early efforts focused on data-driven surrogates that either learn discrete-time evolution maps or infer continuous-time governing equations \cite{qin2019datadriven, raissi2018multistep}, and on sparse-regression schemes that recover governing equations directly from data \cite{brunton2016sindy}. To leverage physical priors and reduce data requirements, Physics-Informed Neural Networks (PINNs) \cite{raissi2019pinn} encode the governing equations as soft constraints in the loss, with extensions to stiff \cite{ji2020stiffpinn}, multi-scale \cite{leung2022nhpinn}, and differential--algebraic systems, including applications in power grids \cite{moya2023daepinn, huang2022pinnpower} and biology \cite{yazdani2020sysbiopinn}. Because PINNs typically require retraining whenever problem parameters change, attention has shifted to operator learning, in which neural architectures approximate mappings between infinite-dimensional function spaces. Deep Operator Networks (DeepONet) \cite{lu2021deeponet, lu2022deeponetcomparison, lu2022multifidelitydeeponet, lin2023responsedeeponet} and Fourier Neural Operators \cite{li2023neuraloperators, wen2022neuraloperator, azizzadenesheli2024neuraloperators} have demonstrated strong generalization across parametric PDE families and complex multi-physics applications, motivating their use as surrogates for safety-critical simulations such as power-system transient analysis \cite{moya2023deeponetgrid}.

\paragraph{Quantum machine learning for scientific computing.}
Parallel to these classical advances, quantum machine learning has progressed toward learning solutions of ordinary and partial differential equations, primarily through variational and hybrid quantum–classical formulations. Variational quantum PDE solvers \cite{joo2021quvapde} and approaches such as quantum lattice-Boltzmann methods augmented with multi-fidelity neural correction \cite{jacob2025multifidelityquantum} extend these ideas to both linear and nonlinear regimes. Quantum Orthogonal Neural Networks (QOrthoNN) \cite{kerenidis2022qorthonn, landman2022medicalqorthonn} replace dense classical layers with parameterized quantum circuits that implement orthogonal transformations, yielding improved asymptotic scaling, and have been empirically tested on medical image classification tasks \cite{landman2022medicalqorthonn}. Building on the QOrthoNN, architectures like Quantum DeepONet \cite{xiao2025qdon} extend operator learning to quantum settings, while independent approaches such as Quantum Fourier Neural Operators (QFNO) \cite{jain2024qfno} have been evaluated on a range of benchmark PDE families. Quantum scientific machine learning has also replicated paradigms such as physics-informed neural networks \cite{trahan2024qpinn, farea2025qcpinn}.

\paragraph{Uncertainty quantification for scientific machine learning.}
Reliable deployment of neural surrogates in safety-critical settings demands rigorous uncertainty quantification \cite{psaros2023uq, zou2024operatoruq}. Classical Bayesian neural networks \cite{mackay1992practical, neal1996bnn} model distributions over parameters but do not scale well or provide mathematical guarantees. Alternative approaches include deep ensembles \cite{lakshminarayanan2017deep, psaros2023uq}, Bayesian PINNs \cite{yang2021bpinn}, probabilistic CNNs for PDEs \cite{winovich2019convpdeuq}, and Bayesian formulations of DeepONet \cite{lin2023bdeeponet}. For a comprehensive review of these deep learning UQ techniques, we refer readers to \cite{abdar2021uqreview} and \cite{psaros2023uq}. These methods generally lack distribution-free coverage guarantees, motivating the adoption of conformal prediction \cite{vovk1999conformal, vovk2022conformallearning, angelopoulous2023conformal} to provide rigorous intervals for neural operators. Recent work further extends conformal methodology to the operator learning paradigm \cite{moya2025conformal}. In the quantum setting, Park and Simeone \cite{park2023qcp} applied conformal prediction to quantum machine learning, but uncertainty-aware quantum operator learning remains largely unexplored.



\section{Problem Formulation}
\label{sec:problemformulation}

We aim to learn an operator $\mathcal{G}$ that maps input functions to pointwise predictions. We train this operator using a supervised dataset $\mathcal{D}=\{(\mathbf{u}_{i},\mathbf{y}_{ij},s_{ij})\}_{i=1,j=1}^{N,M_{i}}$, where $\mathbf{u}_{i}=(u(x_{1}),\dots,u(x_{d_u}))^{\top}\in\mathbb{R}^{d_u}$ represents a discretized input function evaluated at $d_u$ fixed sensor locations. The variable $\mathbf{y}_{ij}\in\mathbb{R}^{d_{y}}$ denotes a spatial or temporal query coordinate, and $s_{ij}=\mathcal{G}(\mathbf{u}_{i})(\mathbf{y}_{ij})\in\mathbb{R}$ represents the corresponding ground-truth value of the operator evaluated at that specific coordinate. To solve this operator learning problem, we use the Deep Operator Network (DeepONet) \cite{lu2021deeponet, lu2022deeponetcomparison, lu2022multifidelitydeeponet, lin2023responsedeeponet} as our foundational approximation architecture.

The DeepONet architecture decomposes the target operator into two sub-networks: a branch network that encodes discretized input functions and a trunk network that embeds spatial or temporal query coordinates. For any generic input function $\mathbf{u}$ and query coordinate $\mathbf{y}$, the network approximates the target operator by computing the inner product of the $p$-dimensional branch output vector, $\mathbf{b}$, and trunk output vector, $\mathbf{t}$:
\[\mathcal{G}(\mathbf{u})(\mathbf{y}) \approx \sum_{k=1}^{p} b_{k}t_{k} = \mathbf{b}^\top\mathbf{t} \in \mathbb{R}\]
We optimize the combined trainable parameters, $\theta \in \mathbb{R}^{d_\theta}$, of both sub-networks simultaneously by minimizing the Mean Squared Error (MSE) across the dataset:
\[\mathcal{L}(\theta) = \frac{1}{N} \sum_{i=1}^{N} \frac{1}{M_{i}} \sum_{j=1}^{M_{i}} (\mathcal{G}_{\theta}(\mathbf{u}_{i})(\mathbf{y}_{ij}) - s_{ij})^{2}\]
where $d_\theta$ is the total number of parameters, and the point-wise predictions and ground truths $s_{ij}$ reside in $\mathbb{R}$. Although this architecture is highly expressive, it introduces a computational bottleneck. Assuming an input size of $d_{in}$ and a uniform hidden layer width of $n$, the initial layer maps the input dimension to $n$ at a cost of $\mathcal{O}(d_{in} \cdot n)$. Subsequent dense hidden layers rely on $\mathcal{O}(n^{2})$. At test time, when the learned operator must be evaluated across thousands of hypothetical conditions in a scenario batch, this quadratic scaling with respect to the network width bounds the total computational budget.

Specifically, offline batch analysis requires evaluating the learned operator across $N$ distinct hypothetical scenarios. To establish an upper bound on computational requirements, let $M = \max_{i}(M_{i})$ denote the maximum number of spatial or temporal query coordinates evaluated per scenario. By caching the branch network's representation, we evaluate it only once per scenario, whereas the trunk network must be evaluated $M$ times. This yields a per-scenario cost of $\mathcal{O}(n^{2} + M \cdot n^{2})$ where $n$ is the size of the hidden layers. Consequently, the total computational cost across the entire batch scales as $\mathcal{O}(N \cdot M \cdot n^{2})$. This explosion in complexity is strictly prohibitive for fine discretizations. Furthermore, any uncertainty quantification method that requires multiple forward evaluations, such as ensembling or Monte Carlo sampling, will linearly multiply this massive baseline cost. Therefore, achieving robust uncertainty quantification at scale requires redesigning the inference mechanism to bypass this classical quadratic bottleneck.

To guarantee reliable uncertainty bounds without exacerbating the quadratic deployment cost, we formalize our objective using conformal prediction. Specifically, for a given input function $\mathbf{u}$ and query coordinate $\mathbf{y}$, we seek to generate adaptive prediction sets, $C_\alpha (\mathbf{u},\mathbf{y})$, that achieve a strict, user-specified miscoverage rate, $\alpha$. Generating these rigorous, distribution-free bounds while remaining computationally feasible for massive scenario batches requires overcoming the classical scaling limits. Simultaneously achieving sub-quadratic operator inference and mathematically calibrated uncertainty defines the core design constraints of our proposed methodology.



\section{Method}
\label{sec:method}

\subsection{Design principles}

\paragraph{Complexity} To eliminate the classical $\mathcal{O}(n^{2})$ inference bottleneck, Quantum Orthogonal Neural Networks (QOrthoNNs) \cite{kerenidis2022qorthonn, landman2022medicalqorthonn} replace dense layers with three sub-circuits: a data loader $S(\cdot)$, a parameterized unitary $U(\cdot)$, and tomography.  Tomography is required to extract the final classical output from the evolved quantum state. Emulating a classical dense layer that executes an orthogonal $m \times n$ linear transformation using a QOrthoNN requires $1+\max(m,n)$ qubits; to establish a direct equivalence with classical hidden layers of width $n$, we evaluate the $n \times n$ case utilizing $n+1$ qubits. The underlying circuit is built using Reconfigurable Beam Splitter (RBS) gates \cite{kerenidis2022qorthonn, landman2022medicalqorthonn, xiao2025qdon}, which require only nearest-neighbour connectivity and admit an $\mathcal{O}(1)$ decomposition into standard hardware basis gates. 

The sequential forward pass begins with the data loader $S(\cdot)$, which, when parameterizing an $n$-dimensional input vector, requires a circuit depth of $n-1$. Subsequently, the pyramidal unitary $U(\cdot)$ executing an $m \times n$ transformation is bounded by a maximum depth of $2\max(m,n)-3$; for our specific $n \times n$ configuration, this simplifies to $2n-3$. Finally, the tomography sub-circuit requires two additional $S(\cdot)$ data loaders alongside $\mathcal{O}(1)$ supplementary gates (X, CNOT, and Hadamard). Cumulatively, this bounds the total circuit depth to $\mathcal{O}(n)$, which can be further compacted by pipelining independent operations. To extract the classical output, the system requires $\mathcal{O}(\log(n)/\delta^{2})$ quantum measurements to bound the finite-sampling error below a threshold $\delta$ \cite{kerenidis2022qorthonn, landman2022medicalqorthonn}. Consequently, the total forward-pass time complexity drops to $\tilde{\mathcal{O}}(n/\delta^{2})$, resolving the classical scalability bottleneck.

\paragraph{Uncertainty Quantification} Although recent work explores leveraging the stochasticity of quantum measurements to quantify uncertainty \cite{park2023qcp}, this methodology is flawed in our case, where epistemic uncertainty cannot be inferred from shot noise. As demonstrated by \cite{xiao2025qdon}, the standard deviation for a single output component $o_j$ of a QOrthoNN scales as:
\[\sigma[o_j] \propto \frac{1}{\sqrt{N_\mathrm{shots}}}\]
Because this standard deviation is proportional to $\tfrac{1}{\sqrt{N_\mathrm{shots}}}$, it vanishes as the measurement budget increases, regardless of input difficulty or out-of-distribution shifts. To capture genuine epistemic uncertainty arising from sparse training data, we must instead construct ensembles of independently trained Quantum DeepONets. However, naïvely executing $L$ independent quantum models scales qubit requirements linearly, introducing a prohibitive hardware bottleneck.

\paragraph{Ensembling} Executing an ensemble of $L$ quantum operators introduces a strict operational trade-off: naïve sequential execution incurs prohibitive temporal overhead due to repeated circuit compilations and state preparations, whereas naïve parallel execution demands an impractical $\mathcal{O}(L \cdot n)$ scaling in spatial qubits. First, a hybrid classical-quantum architecture replaces the low-evaluation-frequency sub-network—such as a branch network evaluated only once per trajectory—with a classical counterpart. For instance, replacing a branch network that is evaluated only once per trajectory with a classical dense network yields a total computational cost of $\mathcal{O}(N \cdot n^2 + N \cdot M \cdot n \log n)$ where $N$ is the number of input trajectories and $M$ is the number of query coordinates per trajectory. Given $M \gg n$, the $\mathcal{O}(n^2)$ classical penalty is safely amortized across the batch. This targeted substitution eliminates the dominant quantum noise contribution of the low-evaluation-frequency sub-network entirely, while rigorously preserving the $\tilde{\mathcal{O}}(n)$ quantum speedup for the high-frequency trunk network evaluations. Second, to efficiently scale the ensemble itself, we deploy Superposed Parameterized Quantum Circuits (SPQCs) \cite{patapovich2025spqc}. By encoding all $L$ ensemble members into a unified circuit via uniformly controlled rotations conditioned on address qubits, SPQCs evaluate the entire ensemble simultaneously. This architecture requires only $n + 1 + \lceil\log_2(L)\rceil$ qubits per sub-network at the cost of linear $\mathcal{O}(L \cdot n)$ circuit depth scaling \cite{patapovich2025spqc}. Guided by these design constraints, we now detail the specific quantum architecture, training procedure, and ensembling methodologies.

\subsection{Quantum DeepONet Architecture}

Each Quantum DeepONet consists of two sub-networks: a branch Quantum Orthogonal Neural Network (QOrthoNN) that processes the input function $\mathbf{u}$, and a trunk QOrthoNN that processes the query coordinate $\mathbf{y}$. Within these sub-networks, Parameterized Quantum Circuits (PQCs) execute orthogonal linear transformations, which are subsequently passed through classical non-linear activations. To remove the strict orthogonality restrictions imposed by unitary quantum operations and increase the model's expressive capacity, a final unconstrained classical linear layer caps each QOrthoNN. The fundamental hardware building block enabling these quantum transformations is the Reconfigurable Beam Splitter (RBS) gate, defined by the unitary matrix:
\[
    U_\mathrm{RBS}(\theta) = 
    \begin{pmatrix}
    1 & 0 & 0 & 0 \\
    0 & \cos\theta & \sin\theta & 0 \\
    0 & -\sin\theta & \cos\theta & 0 \\
    0 & 0 & 0 & 1
    \end{pmatrix}
\]
When applied to a general two-qubit state $\alpha |00\rangle + \beta|01\rangle + \gamma|10\rangle + \delta |11\rangle$, the RBS gate modifies the amplitudes of only the unary $|01\rangle$ and $|10\rangle$ components, mapping the state to $\alpha |00\rangle + (\beta\cos\theta + \gamma\sin\theta)|01\rangle + (-\beta\sin\theta + \gamma\cos\theta)|10\rangle + \delta |11\rangle$. Because it never maps states with zero or two ones into states with a single one, this preservation mechanism allows the entire QOrthoNN to operate strictly within the unary subspace. Consequently, the tomography procedure for extracting the classical output assumes a superposition of unary states, which supports error mitigation techniques.

To clarify the mechanics of a single QOrthoNN hidden layer, we briefly detail the evolution of the quantum state for an $n \times n$ transformation, deferring the generalized $m \times n$ mathematical formulation to Appendix \ref{sec:quantum_state_evolution}. To encode classical data onto the quantum circuit, QOrthoNNs use unary amplitude encoding. Because quantum amplitude encoding requires the input vector to have a strict unity $L_2$-norm, raw classical data must be pre-processed. Raw input vectors undergo a fixed affine min-max normalization into $[-1, 1]$ learned from the training set. To enforce the unity norm constraint required for amplitude encoding, the normalized vector is scaled by $1/\sqrt{n}$, and a slack variable, $\sqrt{1 - \sum (x_i/\sqrt{n})^2}$, is appended. For subsequent hidden layers, intermediate outputs are divided by their $L_2$ norm to ensure unit norm.

Once the $n$-dimensional normalized input vector $\mathbf{x} \in \mathbb{R}^n$ is prepared, the PQC executes the orthogonal linear transformation $\mathbf{W} \in \mathbb{R}^{n \times n}$ using $n$ data qubits and one ancillary qubit. First, a data-dependent unitary $S(\mathbf{x})$ loads the input vector onto the data register via amplitude encoding, yielding the quantum state:
\[
|\psi\rangle_\mathrm{data} = x_1 |10\cdots0\rangle + x_2 |01\cdots0\rangle + \cdots + x_n |00\cdots1\rangle = \sum_{i=1}^{n} x_i |\mathbf{e_i}\rangle
\]
where $\mathbf{e_i}$ denotes a unary computational basis state. This is followed by a pyramidal parameterized unitary $U(\boldsymbol{\theta})$ that emulates the weight matrix $\mathbf{W}$, evolving the data register to:
\[
|\psi\rangle = \sum_{j=1}^{n}\sum_{i=1}^{n}W_{ji}x_i|\mathbf{e_j}\rangle
\]
This operation effectively performs a parallel matrix-vector multiplication directly on the quantum amplitudes. Finally, tomography recovers these transformed amplitudes back into classical floating-point values. To achieve this, the tomography sub-circuit utilizes the ancilla qubit alongside two additional $S$ data loaders. The details of the tomography sub-circuit and quantum state evolution are outlined in Appendix \ref{sec:quantum_state_evolution}. This extracted classical output is then passed through standard non-linear activations to complete the hidden layer sequence. A QOrthoNN is a stack of these ``quantum layers'' followed by a final classical layer, emulating a standard neural network.

Having established the forward-pass mechanics, we must address optimization. To ensure stable convergence and avoid hardware decoherence, we bypass training directly on Noisy Intermediate-Scale Quantum (NISQ) hardware \cite{bharti2022nisq, preskill2018nisq}. The feasibility of classical training is enabled from the fact that these quantum operations mathematically equate to an exact classical orthogonal linear transformation, $\mathbf{W}$. Because this transformation is intrinsically parameterized by the Reconfigurable Beam Splitter (RBS) rotation angles, $\theta$, we can natively emulate the network's exact mathematical behaviour within standard deep learning frameworks, circumventing the need for expensive quantum state evolution simulations. During training, standard automatic differentiation directly optimizes the RBS rotation angles instead of the weight matrix elements in $\mathcal{O}(n^2)$ time \cite{kerenidis2022qorthonn, landman2022medicalqorthonn}. By navigating the optimization trajectory through this angular parameter space, the network guarantees strict mathematical orthogonality by construction. This circumvents the computationally prohibitive $\mathcal{O}(n^{3})$ Singular Value Decomposition (SVD) or Stiefel manifold projection steps required to maintain orthogonality in traditional classical networks \cite{kerenidis2022qorthonn, landman2022medicalqorthonn, jia2019orthogodnn}. Once the angles are optimized classically, these learned parameters are deployed to the PQC to execute the sub-quadratic inference.

This strict mathematical orthogonality provides two critical optimization advantages. First, by preserving eigenvalue magnitudes across deep transformations, the orthogonal constraint prevents the vanishing and exploding gradients that can destabilize deep architectures, thereby promoting robust generalization \cite{kerenidis2022qorthonn, landman2022medicalqorthonn}. Second, within the quantum domain, this orthogonality reshapes the loss landscape. By restricting the optimization trajectory with the orthogonality constraint rather than exploring the exponentially large full Hilbert space, the architecture systematically circumvents the barren plateaus of general PQCs \cite{wang2021barren, mcclean2018barren, schumann2024barren}.

\subsection{Ensemble Construction}

To induce predictive diversity, we train each ensemble member, $\mathcal{G}_{\theta_m}$, independently on the full dataset using distinct random shuffling and parameter initializations. At inference, the resulting ensemble mean, $\mu(\mathbf{u})(\mathbf{y}) = \frac{1}{L} \sum_{m=1}^L \mathcal{G}_{\theta_m}(\mathbf{u})(\mathbf{y})$, provides the central prediction, while the ensemble standard deviation, $\sigma(\mathbf{u})(\mathbf{y})$, quantifies the epistemic disagreement. Although we evaluated non-parametric bootstrapping (bagging) as an alternative mechanism, we observed no performance gain. Consistent with classical deep learning theory \cite{lakshminarayanan2017deep, lee2015mheads}, because these high-capacity models require maximal training data, the loss of unique samples inherent to bootstrapping outweighs any diversity benefits. Finally, because sequentially executing this identically distributed ensemble remains prohibitive on near-term hardware, we address resource efficiency through two specialized execution strategies.

\subsection{Hybrid Classical–Quantum Architectures}

The DeepONet's inherent branch-trunk decomposition allows us to selectively replace one sub-network with a classical counterpart to optimize resource efficiency. Specifically, we replace the low-evaluation-frequency component—typically the branch network, which is evaluated only once per input trajectory—to eliminate its quantum hardware noise penalty. This classical substitution incurs a one-time $\mathcal{O}(n^2)$ computational penalty per trajectory. However, because the total evaluation cost scales as $\mathcal{O}(N \cdot n^2 + N \cdot M \cdot n \log n)$, this substitution remains optimal whenever the number of query coordinates per trajectory significantly exceeds the network width ($M \gg n$). Under this condition, the $\mathcal{O}(n^2)$ classical penalty is safely amortized, successfully isolating and eliminating the dominant hardware noise contribution while rigorously preserving the $\tilde{\mathcal{O}}(n)$ quantum acceleration for the high-frequency trunk network.

Alternatively, one could reduce resource demands by sharing a single sub-network across all ensemble members. However, this strategy sacrifices predictive diversity. By forcing the independent, unshared sub-networks to adapt to a fixed intermediate output, this shared representation acts as an informational bottleneck. Thus, it artificially constrains the ensemble's variance and degrades uncertainty quantification by yielding overconfident predictions.

\subsection{Superposed Parameterized Quantum Circuits}

To eliminate the linear qubit scaling associated with parallel execution, our second strategy employs Superposed Parameterized Quantum Circuits (SPQCs) \cite{patapovich2025spqc}. By efficiently encoding all $L$ ensemble members into a single unified circuit, SPQCs require only $n + 1 + \lceil\log_2(L)\rceil$ qubits per sub-network for $n \times n$ transformations at the cost of linear $\mathcal{O}(L \cdot n)$ circuit depth scaling. We realize this architecture through a cascade of uniformly controlled rotations conditioned on address qubits initialized in a uniform superposition:
\[
    |\psi\rangle = \frac{1}{\sqrt{L}} \sum_{j=1}^{L}
    |j\rangle_{\mathrm{addr}} \otimes U(\boldsymbol{\theta}^{(j)})\,S(\mathbf{x}^{(j)})\,|\mathbf{e}_1\rangle_{D},
\]
Here $|\mathbf{e}_1\rangle_{D}$ is the basis state with the first data qubit set to $|1\rangle$ and all others $|0\rangle$, in which each data register is initialized before encoding (Appendix \ref{sec:quantum_state_evolution}). By conditioning the data loading unitaries, $S(\mathbf{x}^{(j)})$, on these address bits, we ensure each specific model, $U(\boldsymbol{\theta}^{(j)})$, processes its correct input data. Correspondingly, a modified tomography protocol extracts the independent outputs of all $L$ models simultaneously. This multiplexing enables the system to perform state preparation, unitary evolution, and measurement exactly once for the entire ensemble. Although the requisite controlled operations inevitably increase circuit depth and exacerbate accumulated quantum noise on near-term devices, SPQCs provide a highly scalable, resource-efficient execution strategy as physical hardware coherence improves. 

While executing models sequentially requires swapping $L$ distinct circuits—introducing significant temporal loading overhead—and naïvely deploying $L$ parallel circuits demands a prohibitive spatial scaling of $\mathcal{O}(L \cdot n)$ qubits, the SPQC encodes the entire ensemble into a unified circuit exactly once. By circumventing these execution bottlenecks, this multiplexed approach makes hardware utilization more efficient. However, this introduces a trade-off: to maintain fixed per-model precision, the total shot budget must still scale linearly with the ensemble size $L$, as the measurement samples are probabilistically distributed across the address branches.




\section{Uncertainty Quantification}
\label{sec:uq}

To translate the raw ensemble dispersion into mathematically calibrated bounds, we apply adaptive conformal prediction \cite{angelopoulous2023conformal, vovk1999conformal, vovk2022conformallearning, shafer2008conformaltutorial, papadopoulos2002conformalregression}. During the calibration phase, we compute an adaptive nonconformity score, $r_{ij}$, for each data point $(\mathbf{u}_i, \mathbf{y}_{ij}, s_{ij})$ by normalizing the absolute prediction error by the local ensemble standard deviation:
\[r_{ij} = \frac{|s_{ij} - \mu(\mathbf{u}_i)(\mathbf{y}_{ij})|}{\sigma(\mathbf{u}_i)(\mathbf{y}_{ij}) + \epsilon}\]
By calculating the appropriate empirical quantile, $\hat{q}$, of these nonconformity scores, we construct an adaptive uncertainty bound around the predicted ensemble mean. Specifically, for $N_{cal}$ exchangeable calibration examples with sorted scores $r_{(1)} \le \dots \le r_{(N_{cal})}$, we define the quantile as $\hat{q} = r_{(k)}$ where $k = \lceil(N_{cal}+1)(1-\alpha)\rceil$, with $\hat{q} = +\infty$ if $k > N_{cal}$. For a given input function $\mathbf{u}$ and a query coordinate $\mathbf{y}$, any candidate true output value $v \in \mathbb{R}$, the prediction set is defined as:
\[C_{\alpha}(\mathbf{u}, \mathbf{y}) = \{v \in \mathbb{R} : |v - \mu(\mathbf{u})(\mathbf{y})| \le \hat{q} \cdot (\sigma(\mathbf{u})(\mathbf{y}) + \epsilon)\}\]
Crucially, when the calibration and test examples remain exchangeable, this framework provides a distribution-free mathematical guarantee of finite-sample pointwise marginal coverage \cite{angelopoulous2023conformal, shafer2008conformaltutorial}. This ensures that at any specific query coordinate, the true output value falls within the prediction set with a strict, user-specified probability:
\[P(s_{ij} \in C_{\alpha}(\mathbf{u}_i, \mathbf{y}_{ij})) \ge 1 - \alpha\]

The theoretical validity of these distribution-free coverage bounds requires the calibration and test examples to be exchangeable \cite{angelopoulous2023conformal, shafer2008conformaltutorial}. However, real-world deployments introduce challenges to this assumption, which we explicitly acknowledge. First, physical quantum hardware noise—such as fluctuating depolarizing channels, readout errors, and thermal relaxation—can drift between the calibration and test phases, inherently threatening exchangeability. While we currently evaluate conformal prediction assuming these noise parameters remain static between calibration and testing, future deployments on temporally drifting hardware will formally violate this premise. Second, evaluating sequential tasks with temporal dependencies, such as online prediction with sliding windows, violates this statistical assumption by definition. Rather than restricting our analysis to idealized conditions, we treat the non-exchangeable online setting as a rigorous stress test for our uncertainty calibration. We evaluate the framework's empirical coverage under both exchangeable batch settings (with static noise) and this non-exchangeable sequential environment, deferring the derivation of formal theoretical bounds for the latter to future work.

To clarify our calibration protocol, we explicitly define our sampling units. For online prediction with sliding windows, the temporal data structure restricts the sampling unit to a single memory window paired with its immediate next-step prediction. For other offline batch settings, the fundamental unit of exchangeability is the entire input trajectory. However, to construct the empirical nonconformity distribution, we pool the individual scores calculated at every evaluated query coordinate across all calibration trajectories, rather than computing a single aggregated score per trajectory.



\section{Experiments}
\label{sec:experiments}

\subsection{Setup and Evaluation Metrics}

We evaluate the predictive accuracy of our framework using relative $L_2$ error, and we assess the quality of uncertainty calibration through empirical coverage, average interval width, and peak uncertainty across both synthetic and real-world operator-learning tasks. Establishing an ideal baseline for these metrics on large quantum circuits, those exceeding 20 qubits, presents a severe computational bottleneck, as traditional state-vector simulations incur exponential memory costs. To bypass this limitation, we execute ideal simulations by directly evaluating the mathematically equivalent classical Orthogonal Neural Network, perfectly mimicking the noiseless quantum state evolution while avoiding exponential overhead.

Beyond these ideal baselines, we evaluate the framework under varying degrees of quantum noise. To conduct hyperparameter ablations efficiently, we first use a simplified gate-based noise model incorporating finite-sampling error and depolarizing channels. For single- and two-qubit gates acting on a $k$-qubit subsystem $A$ (where $k \in \{1, 2\}$), the local depolarizing error on the quantum state $\rho$ is modeled as:
\[\mathcal{E}_A(\rho) = (1-\lambda)\rho + \lambda \left( \frac{I_A}{2^k} \otimes \text{Tr}_A[\rho] \right)\]
with noise strength $\lambda$ scaled by 0.8 for two-qubit operations to equalize error rates across 1 and 2-qubit depolarizing noise channels. We evaluate circuits under this model by transpiling them to IBM's Eagle basis gate set $\{\mathrm{ECR}, \mathrm{R}_z, \mathrm{SX}, \mathrm{X}, \mathrm{I}\}$ \cite{abughanem2025eagle, xiao2025qdon}, assuming nearest-neighbour connectivity. To accelerate simulations, we perform multinomial sampling of the evolved density matrix, which perfectly correlates with computationally expensive per-shot simulations as outlined in Appendix \ref{sec:multinomial}. We rigorously validate near-term hardware readiness using a Qiskit Aer noise model \cite{qiskit2024}. This realistic model simulates single-qubit readout errors, depolarizing errors, and thermal relaxation. For this validation, circuits are strictly transpiled according to the reported calibration parameters, basis gates, and topological coupling maps of three representative quantum processing units: IBM Brisbane (127-qubit Eagle r3), IBM Torino (133-qubit Heron r1), and IBM Marrakesh (156-qubit Heron r2) \cite{IBMQuantumComputers}. The snapshots of calibration data used in our experiments can be found in Appendix \ref{sec:snapshots}.

Operating in these realistic, noisy environments needs robust error mitigation. Because the QOrthoNN architecture structurally encodes data exclusively within the unary subspace, the valid quantum state must always maintain a Hamming weight of exactly one. We exploit this strict physical constraint to mitigate hardware errors through post-selection. During tomography, we simply discard any measurement shot that collapses into an invalid, non-unary state. By trading a larger total shot budget for higher fidelity readouts, this post-selection systematically filters out corruption, reducing the impact of hardware noise on the final operator predictions.


\subsection{Synthetic Benchmarks}

\begin{table}[t]
\caption{Summary of model performance across synthetic experiments under ideal simulation with 90\% target coverage. The metrics from left to right are Relative $L_2$ Error (\%), Coverage (\%), Average Width, and Peak Uncertainty.}
\label{tab:results_summary_syn}
\begin{center}
\begin{tabular}{lcccc}
\hline
\textbf{Experiment} & \textbf{Rel. $L_2$ Err. (\%)} & \textbf{Cov. (\%)}  & \textbf{Avg. Width} & \textbf{Peak Unc.} \\
\hline

Antiderivative (Ens. Size = 4) & 0.46 & 88.40 & 0.004 & 0.044 \\
Antiderivative (Ens. Size = 8) & 0.46 & 92.13 & 0.005 & 0.080 \\
Advection (Ens. Size = 4)      & 2.38 & 89.36 & 0.062 & 0.751 \\
Advection (Ens. Size = 8)      & 2.28 & 88.99 & 0.053 & 0.621 \\

\hline
\end{tabular}
\end{center}
\end{table}

\begin{table}[t]
\caption{Architecture and hardware specifications for the quantum DeepONet ensembles for the synthetic benchmarks. Width denotes the number of neurons per layer. Params and circuit depth report the maximum per-layer values for total RBS parameters and transpiled PQC depth (on IBM Heron/Eagle), respectively. All models use the SiLU activation function. (Config: R = ResNet, F = Fourier features). The number of qubits used for a single model is $d_u + 2$ (accounting for the $d_u$ raw inputs, one appended slack coordinate, and one measurement ancilla).}
\label{tab:model_architectures_syn}
\begin{center}
\begin{tabular}{lccccccccc}
\hline
\textbf{Experiment} & $d_u$ & $d_y$ & \textbf{Layers} & \textbf{Width} & \textbf{Params} & \textbf{Heron} & \textbf{Eagle} & \textbf{Config} & \textbf{Ens. Size} \\
\hline 

Antiderivative & 10 & 1 & 2 & 10 & 55  & 165 & 321 & -- & 4/8 \\
Advection      & 20 & 1 & 7 & 20 & 210 & 305 & 601 & R  & 4/8 \\

 \hline
\end{tabular}
\end{center}
\end{table}

\begin{figure}[t]
\begin{center}
\includegraphics[width=\columnwidth]{figures/antiderivative_noise_analysis.pdf}
\end{center}

\caption{Impact of depolarizing noise $\lambda$ (represented by different colours) and finite-sampling error (x-axis) on Antiderivative prediction performance. All three metrics: $L_2$ Error (left), Average Width (middle), and Max Width (right) improve with increasing shot counts and yield better performance at lower noise levels.}

\label{fig:simplified_anti}
\end{figure}

\subsubsection{Antiderivative Operator}

We first evaluate our framework on a synthetic antiderivative operator, $\mathcal{F}:v \mapsto u$. This operator maps an input function $v(x)$ to its integral $u(x)$ over the domain $x\in[0,1]$, defined by the differential relation $\frac{du(x)}{dx}=v(x)$ subject to the initial condition $u(0)=0$. To generate the dataset, we sample smooth input functions from a Gaussian Random Field (GRF) prior, $v\sim\mathcal{GP}(0,k_l(x_1,x_2))$. The covariance is governed by a squared exponential kernel, $k_l(x_i,x_j)=\exp(-\frac{d(x_i,x_j)^2}{2l^2})$, where the length scale $l$ controls the smoothness of the sampled functions. For this task, we deploy an ensemble of $L=8$ independent models, using a 2-layer, 10-neuron QOrthoNN architecture. As detailed in Table \ref{tab:model_architectures_syn}, this model requires only 55 parameters per layer and transpiles to a maximum circuit depth of 165 on the IBM Heron topology and 321 on IBM Eagle. Under ideal noiseless simulation, this quantum ensemble successfully learns the operator, achieving a 0.46\% relative $L_2$ error and 92.13\% empirical coverage for a target rate of $1-\alpha=0.90$. 

We further studied the effects of depolarizing noise on the approximation of the antiderivative operator with $L=8$ ensemble members each composed of the same QOrthoNN architecture described above. We observed that accuracy and uncertainty metrics follow a general expected trend of better performance with decreasing depolarizing noise levels and finite-sampling error (increasing shot count) as can be observed in Figure \ref{fig:simplified_anti}.


\subsubsection{Advection Equation}

To evaluate the framework on a dynamical system, we next consider the 1D advection equation, $\frac{\partial u}{\partial x} + \frac{\partial u}{\partial t} = 0$ for $x, t \in [0, 1]$, subject to periodic boundary conditions. Here, the operator learns the solution mapping from the initial condition $u(x,0)$ to the full spatiotemporal trajectory $u(x,t)$. To enforce the required periodicity, we sample the initial conditions from a Gaussian Random Field utilizing an Exp-Sine-Squared kernel: $k_l(x_i, x_j) = \exp \left( \frac{-2\sin^2(\pi d(x_i, x_j)/p)}{l^2} \right)$. Because advection demands tracking features across spacetime, we deploy a deeper, 7-layer QOrthoNN architecture featuring residual connections to stabilize gradient flow. This model processes 20 uniformly sampled spatial locations for the branch input and a Cartesian product of 50 spatial and temporal coordinates for the trunk query. Consequently, this expanded architecture requires 210 parameters per layer, transpiling to a maximum circuit depth of 305 on IBM Heron and 601 on IBM Eagle. Under ideal simulation, the 8-model advection ensemble successfully maps the dynamics, achieving a 2.28\% relative $L_2$ error and 88.99\% empirical coverage for a target rate of $1-\alpha=0.90$. Table~\ref{tab:results_summary_syn} summarizes the accuracy and calibration metrics for both ensemble sizes across the synthetic benchmarks. Having established these theoretical baselines, we now evaluate whether these rigorous conformal coverage guarantees survive under realistic, hardware-calibrated quantum noise.




\subsection{Coverage Under Realistic Noise}

\begin{figure}[t]
\begin{center}
\includegraphics[width=\columnwidth]{figures/cov_guarantee.pdf}
\end{center}

\caption{Empirical coverage of the prediction intervals on the test set for the antiderivative task using the complex noise model. We simulated noise models calibrated to IBM Marrakesh, IBM Torino, and IBM Brisbane. The dashed line indicates the target nominal coverage level of 90\%. The results demonstrate that the conformal framework maintains validity ($\ge 90\%$) even under the more realistic noise model.}

\label{fig:cov_guar}
\end{figure}

To validate practical robustness, we evaluate the empirical coverage of the antiderivative ensemble under realistic noise models calibrated to the IBM Marrakesh, Torino, and Brisbane quantum processors. As illustrated in Figure \ref{fig:cov_guar}, when evaluated across measurement budgets ranging from 5,000 to 100,000 shots using full per-shot simulations rather than multinomial sampling, all three hardware noise profiles successfully maintain empirical coverage at or above the target $1-\alpha=0.90$ threshold. Crucially, this confirms that as long as the hardware error profile—encompassing depolarizing errors, thermal relaxation, and single-qubit readout errors—remains stationary between the calibration and test phases, the inherently noisy quantum execution preserves the exchangeability assumption required for valid conformal prediction. Having demonstrated both sub-quadratic inference and theoretical robustness under physical hardware constraints, we now apply the framework to complex, real-world dynamical systems.



\subsection{Real Power-Engineering Datasets}

Transitioning to real-world power engineering, the complex dynamics of these physical systems impose distinct operational requirements that require specializing our framework into three distinct operator classes. The first is the Transient-to-Transient Operator ($\mathcal{F}$), which maps an entire temporal input trajectory directly to an entire output trajectory. Because this global mapping evaluates the complete dynamic evolution simultaneously, it is essential for offline applications such as post-mortem event analysis and long-term system planning \cite{kwatny2016psdynamics, kundur2004psdynamics}. Formally, this operator functions across the Hilbert spaces $\mathcal{U} = L^2([0,T]; \mathbb{R}^d)$ and $\mathcal{Y} = L^2([0,T]; \mathbb{R}^d)$, where $d$ represents the state-space dimensionality. Utilizing a dataset of observed input-output trajectory pairs, $\{(\mathbf{u}_i, \mathbf{s}_i)\}_{i=1}^N$ where $\mathbf{s}_i = \mathcal{F}[\mathbf{u}_i]$, the objective is to learn a highly accurate quantum surrogate, $\hat{\mathcal{F}}$. We then construct the corresponding pointwise conformal prediction sets, $C_\alpha(\mathbf{u}_i, t)$ for all $t \in [0, T]$, to guarantee reliable inference over the entire simulation horizon.

The second class is the Causal Forecasting Operator $(\mathcal H)$, which supports proactive grid management strategies such as long-range dynamic stability assessments. Unlike the global mapping of the transient operator, $\mathcal H$ explicitly enforces causality by mapping past pre-event histories to future post-event predictions. Formally defined as $\mathcal H: \mathcal U_\mathrm{past} \to \mathcal Y_\mathrm{future}$, this operator processes historical input trajectories from the space $\mathcal{U}_{\text{past}} = L^2([t_0 - \tau, t_0]; \mathbb{R}^d)$ to predict the system's future dynamic evolution within the space $\mathcal{Y}_{\text{future}} = L^2([t_0, t_0 + T_f]; \mathbb{R}^d)$, where $t_0$ denotes the time of a critical grid contingency and $d$ represents the physical state-space dimensionality.

The final class is the Causal Point-wise Transient Operator $(\mathcal G)$, designed for real-time applications such as autonomous grid monitoring and closed-loop control. Mathematically, this is formulated as a time-indexed family of operators, $\{\mathcal{G}_t\}_{t \in [0,T]}$. Each specific operator, $\mathcal{G}_t : L^2([0,t]; \mathbb{R}^d) \to \mathbb{R}^d$, maps a historical input trajectory restricted to the current time $[0,t]$ to a single, localized state prediction at the immediate future step, $t + \Delta t$. However, because integrating a continuously expanding historical domain from $t=0$ is computationally prohibitive for high-frequency real-time inference, we approximate this mapping using a fixed memory window, $\tau > 0$. This ensures operational tractability by truncating the evaluation to a sliding window: $\mathcal{G}_t[u] \approx \mathcal{M}_t\left[u(s)\big|_{s = t-\tau}^t\right]$. Having formally defined these three specialized operator classes, we now evaluate the framework across a series of highly complex, real-world power engineering tasks, with the resulting accuracy and calibration metrics summarized in Table~\ref{tab:results_summary_all}.

\begin{table}[t]

\caption{Summary of model performance across online and offline experiments under ideal simulation. The target coverage for all uncertainty quantification tasks is 90\%.}
\label{tab:results_summary_all}

\begin{center}
\begin{tabular}{lcccc}
\hline
\textbf{Experiment} & \textbf{Rel. $L_2$ Error (\%)} & \textbf{Coverage (\%)} & \textbf{Avg. Width} & \textbf{Peak Uncertainty} \\
\hline

Online V-to-V  & 5.49  & 90.11 & 0.148 & 1.328 \\
Offline V-to-V & 12.49 & 89.66 & 0.345 & 2.691 \\
Offline V-to-P & 4.08  & 89.74 & 0.001 & 0.063 \\

\hline
\end{tabular}
\end{center}

\end{table}

\begin{table}[t]

\caption{Architecture and hardware specifications for the quantum DeepONet ensembles across online and offline tasks. Width denotes the number of neurons per layer. Params and circuit depth report the maximum per-layer values for total RBS parameters and transpiled PQC depth (on IBM Heron/Eagle), respectively. All models use the SiLU activation function. (Config: R = ResNet, F = Fourier features). The number of qubits used for a single model is $d_u + 2$ (accounting for the $d_u$ raw inputs, one appended slack coordinate, and one measurement ancilla).}
\label{tab:model_architectures_all}

\begin{center}
\begin{tabular}{lccccccccc}
\hline
\textbf{Experiment} & $d_u$ & $d_y$ & \textbf{Layers} & \textbf{Width} & \textbf{Params} & \textbf{Heron} & \textbf{Eagle} & \textbf{Config} & \textbf{Ens. Size} \\
\hline

Online V-to-V  & 5  & 1 & 2 & 5  & 15  & 94  & 181 & --   & 4 \\
Offline V-to-V & 20 & 1 & 6 & 20 & 210 & 305 & 601 & R, F & 8 \\
Offline V-to-P & 20 & 1 & 6 & 20 & 210 & 305 & 601 & R, F & 8 \\

\hline
\end{tabular}
\end{center}

\end{table}

\subsection{Offline voltage-to-voltage}

We test the Causal Forecasting Operator $(\mathcal H)$ on an offline voltage prediction task, mapping a 1.9-second pre-fault history to a 2.0-second post-fault trajectory. The input history is discretized into a branch vector $u\in \mathbb R^{20}$, while the forecast horizon is queried across 100 uniformly spaced trunk coordinates. To mitigate spectral bias caused by severe fault transients, we augment the scalar temporal query, $t$, with Fourier features. We extract the $K=5$ dominant frequencies via FFT with a Hanning window, mapping $t$ to an 11-dimensional feature vector:
\[
    \mathbf{y}(t) = [t, \cos(2\pi f_1 t), \sin(2\pi f_1 t), \dots, \cos(2\pi f_5 t), \sin(2\pi f_5 t)]^{\top} \in \mathbb{R}^{11}.
\]

To process these signals, we deploy an $L=8$ ensemble of 6-layer, 20-neuron QOrthoNNs with residual connections. This configuration requires 210 parameters per layer, transpiling to a maximum quantum circuit depth of 305 on IBM Heron and 601 on IBM Eagle. Under ideal simulation, this framework achieves a 12.49\% relative $L_2$ error. Furthermore, the conformal mechanism yields an 89.66\% empirical coverage, effectively maintaining the target threshold of $1-\alpha=0.90$ despite the highly oscillatory dataset.

\subsection{Offline voltage-to-active power}

We next evaluate the Transient-to-Transient Operator $(\mathcal F)$ on an offline cross-domain task, mapping a 10-second voltage trajectory to its corresponding active power trajectory. The voltage input is discretized into a branch vector $v\in \mathbb R^{20}$, while the active power output is queried across 100 Fourier-augmented trunk coordinates. To ensure quantum simulation feasibility against extreme high-frequency transients, the target active power trajectories are smoothed via a zero-phase forward-backward fourth-order Butterworth low-pass filter with a cutoff frequency of $f_c = f_s / 100$. Thus, the network learns to approximate the filtered operator, formally defined as $\mathcal{F}: v \mapsto p_\mathrm{filtered}$.

This task deploys the identical $L=8$ ensemble architecture used in the prior forecasting benchmark, featuring 6-layer, 20-neuron QOrthoNNs with equivalent parameter counts and quantum circuit depths. Under ideal simulation, this ensemble achieves a 4.08\% relative $L_2$ error. The conformal mechanism again proves robust, delivering an 89.74\% empirical coverage against the strict $1-\alpha=0.90$ target threshold. 

\subsection{Online voltage-to-voltage}

\begin{figure}[t]

\begin{center}
\includegraphics[width=\columnwidth]{figures/online_voltage_noise_analysis.pdf}
\end{center}

\caption{Impact of depolarizing noise $\lambda$ on Online Voltage-to-Voltage prediction performance (shots$=10^5$). Mean relative $L_2$ error improves at higher noise levels (left). Empirical coverage also improves at higher noise levels (middle). Evolution of average and maximum prediction interval widths (right).}

\label{fig:online_analysis}
\end{figure}

\begin{figure}[t]

\begin{center}
\includegraphics[width=\columnwidth]{figures/online_pred.pdf}
\end{center}

\caption{Time-varying prediction intervals constructed by conformal prediction on the online voltage prediction task for a high-variance signal portrayed in different colours to differentiate distinct noise levels. The ground truth is drawn in black and the mean prediction is a solid line of the corresponding colour. Figure (b) shares the same y-axis as figure (a). Interestingly, the average and max width (peak uncertainty) decrease as noise level increases for this particular sample.}

\label{fig:online_pred}
\end{figure}

To conclude the real-world evaluations, our final benchmark tests the Causal Point-wise Transient Operator $(\mathcal G)$ under a non-exchangeable data structure. We evaluate the Causal Point-wise Transient Operator $(\mathcal G)$ on an online voltage prediction task using a sliding memory window of $\tau=10$. For this real-time application, we deploy a compact $L=4$ ensemble of 2-layer, 5-neuron QOrthoNNs. This architecture requires only 15 parameters per layer, transpiling to a maximum circuit depth of 94 on IBM Heron. Under ideal simulation, the ensemble achieves a 5.49\% relative $L_2$ error and 90.11\% empirical coverage against the $1-\alpha=0.90$ target. Crucially, because the sliding-window data structure introduces strict temporal dependencies, the conformal exchangeability assumption is inherently violated. Thus, this 90.11\% result represents a purely empirical observation rather than a distribution-free mathematical guarantee.

Because simulating quantum noise across massive, sequential datasets is computationally intractable, we evaluate the framework's robustness by applying the simplified depolarizing noise model to targeted data subsamples, including a high-variance, single-signal test set. As illustrated by metrics in Figure \ref{fig:online_analysis} and the prediction intervals in Figure \ref{fig:online_pred}, these ablations reveal a counterintuitive scaling effect. For shallow circuits using 7 or fewer qubits, desirable uncertainty quantification properties, specifically empirical coverage and average interval width, actually improved at higher depolarizing noise levels $(\lambda)$. We discuss the mechanisms driving this anomalous, noise-induced regularization in the limitations section. 

It must be noted that the training and evaluation procedures differed between these two figures, leading to distinct manifestations of this effect. For the aggregated metrics (Figure \ref{fig:online_analysis}), the models were trained and evaluated on randomly subsampled windows drawn from the full signal dataset. In this generalized regime, increased noise improved overall coverage, predictive accuracy, and maximum interval width. Conversely, the visualization (Figure \ref{fig:online_pred}) isolates a single high-variance test signal. To satisfy the distributional matching required for conformal prediction to function in this specific qualitative example, the corresponding training set was restricted to high-variance windows. Under these conditions, increased noise primarily drove a significant reduction in the average interval width. The varying expressions of this noise-induced advantage between the two figures are artefacts of these contrasting training and inference procedures.

\subsection{Hybrid architecture}

\begin{figure}[t]

\begin{center}
\includegraphics[width=\columnwidth]{figures/hybrid_architecture_std_dev_plot.pdf}
\end{center}

\caption{Performance comparison of hybrid classical-quantum architectures on the antiderivative operator task. We evaluate three ensemble configurations: Fully Quantum (purple), Hybrid with Classical Trunk (orange), and Hybrid with Classical Branch (green). The shaded regions indicate $\pm 1$ standard deviation, where the aggregated predictions were evaluated across a noise parameter space formed by the cross-product of hardware noise (depolarizing levels $\lambda \in \{0.0002, 0.0004, 0.0006, 0.0008\}$) and finite sampling (shot counts ranging from $10^3$ to $10^8$).}

\label{fig:hybrid_arch}
\end{figure}

We isolate the dominant hardware noise source by evaluating fully quantum, classical-branch, and classical-trunk configurations on the antiderivative task. For this specific operator, the quantum branch network is significantly more complex than the trunk, causing it to accumulate substantially more hardware noise and act as the primary fidelity bottleneck. As illustrated in Figure \ref{fig:hybrid_arch}, the classical-branch hybrid achieves the lowest relative $L_2$ error and the narrowest prediction bands across the entire simulated noise space by replacing this bottleneck with a stable classical counterpart. In contrast, the classical-trunk hybrid yields only marginal improvements over the fully quantum baseline because it retains the noisy quantum branch. This confirms that optimal hybrid deployment depends on identifying and replacing the specific sub-network contributing the dominant noise penalty.

\subsection{Superposed Parameterized Quantum Circuits}


We evaluate Superposed Parameterized Quantum Circuits (SPQCs) to determine whether compressing an ensemble into a single unified circuit successfully preserves predictive accuracy and uncertainty quantification. Tested on the antiderivative task using an $L=4$ ensemble of 2-layer, 5-neuron QOrthoNNs, the SPQC architecture requires only 8 qubits per sub-network, a strict spatial reduction from the 24 qubits ($L(n+1)$, each ensemble member has its own measurement ancilla) mandated by naïve parallel execution. As illustrated in Figure \ref{fig:SPQC_comparison}, this unified SPQC matches the standard ensembling baseline across all predictive and uncertainty metrics, maintaining calibration even when subjected to varying levels of depolarizing noise.

Table \ref{tab:spqc_vs_normal_gates} details the corresponding gate decomposition transpiled onto the IBM Heron processor. While the required count of $cz$ gates inevitably increases to accommodate the controlled rotations, the counts of the remaining basis gates avoid linear scaling with the ensemble size. Furthermore, pipelining individual RBS gates inherently restrains the maximum SPQC circuit depth to 272. This remains substantially more efficient than the cumulative depth of 380 required to execute four independent, 95-depth standard circuits sequentially. These metrics physically validate the SPQC as a highly scalable, resource-efficient execution strategy whose comparative advantage will compound as hardware coherence improves.

\begin{figure}[t]

\begin{center}
\includegraphics[width=\columnwidth]{figures/SPQC_comparison_plot.pdf}
\end{center}

\caption{Impact of depolarizing noise $\lambda$ on the standard ensembling framework and the SPQC architecture (shots=$10^5$). The SPQC tracks the standard ensembling framework across all metrics.}

\label{fig:SPQC_comparison}
\end{figure}

\begin{table}[t]

\caption{Circuit resource comparison between normal and SPQC execution modes for a single layer. Width denotes the number of neurons per layer. Circuit depth and gate counts are reported for the transpiled circuits.}
\label{tab:spqc_vs_normal_gates}

\begin{center}
\begin{tabular}{lcccccccccc}
\hline
\textbf{Mode} & $d_u$ & $d_y$ & \textbf{Layers} & \textbf{Width} & \textbf{Depth} & \textbf{CZ} & \textbf{RX} & \textbf{RZ} & \textbf{SX} & \textbf{X} \\
\hline

Normal & 5 & 1 & 2 & 5 & 95  & 62  & 76  & 62  & 21 & 1 \\
SPQC   & 5 & 1 & 2 & 5 & 272 & 173 & 144 & 125 & 44 & 1 \\

\hline
\end{tabular}
\end{center}

\end{table}

\section{Limitations and Outlook}
\label{sec:limitations}

Our conformal coverage guarantees rely on data exchangeability. While satisfied in standard batch scenarios, this assumption is inherently violated by temporally dependent streaming data and potentially compromised by correlated quantum noise, such as hardware cross-talk. To safely formalize our empirical online coverage for real-time deployment, future work must extend this framework via adaptive techniques like weighted or sequential conformal prediction. Beyond theoretical calibration, we must physically investigate the anomalous noise-induced regularization observed in shallow circuits to determine if hardware decoherence can be actively harnessed to improve uncertainty bounds in larger models. Finally, subsequent studies must systematically quantify the operational trade-offs between the immediate fidelity advantages of classical-quantum hybrids and the asymptotic resource efficiency of unified SPQCs.

Currently, our empirical validation remains inherently restricted to compact quantum circuits due to the exponential computational overhead of classical simulation. However, as physical hardware coherence inevitably matures, SPQC-based ensembling—driven by its logarithmic spatial scaling—will emerge as the dominant, highly scalable execution strategy. Transitioning this rigorous framework from offline batch analysis to real-time, uncertainty-aware operator learning represents the critical next step toward deploying robust quantum surrogates in complex, safety-critical infrastructure.

\section*{Code Availability}
The code supporting the findings of this study is publicly available in the GitHub repository: \href{https://github.com/purav-0000/conformalized-quantum-deeponet}{https://github.com/purav-0000/conformalized-quantum-deeponet}. The code is released to facilitate reproducibility and further research, and it may be updated over time as improvements are made.

\section*{Acknowledgment}
We would like to thank the support of National Science Foundation (DMS-2533878, DMS-2053746, DMS-2134209, ECCS-2328241, CBET-2347401 and OAC-2311848), and U.S.~Department of Energy (DOE) Office of Science Advanced Scientific Computing Research program DE-SC0023161, the SciDAC LEADS Institute, and DOE–Fusion Energy Science, under grant number: DE-SC0024583.



\bibliographystyle{vancouver}
\bibliography{main}



\appendix
\section*{Appendix}

\section{Detailed QOrthoNN quantum state evolution}
\label{sec:quantum_state_evolution}

For an orthogonal $m \times n$ transformation $\mathbf{W}$ the required number of data qubits is $q = \max(m, n)$. Let $|\mathbf{e}_k\rangle$ denote the physical unary basis state where the $k$-th qubit is in the $|1\rangle$ state.
\begin{enumerate}

    \item \begin{enumerate}
        \item If $m = n$, then $q = n$ and all data qubits are used at every stage of the circuit.
        \item If $m > n$, then $q = m$. The classical data is encoded onto only the bottom $n$ qubits (indices $m - n + 1$ to $m$), and all $m$ qubits are measured.
        \item If $m < n$, then $q = n$. The classical data is encoded on all $n$ qubits (indices $1$ to $n$), and only the bottom $m$ qubits (indices $n - m + 1$ to $n$) are measured.
    \end{enumerate}

    \item \textbf{State preparation} — For an $n$-dimensional input vector $\mathbf{x}$ and an $m \times n$ linear transformation, $q$ qubits are initialized in the ground state along with an ancilla qubit in uniform superposition. The ancilla qubit is used to infer the sign of the components of the output vector since a quantum measurement natively yields positive probabilities. 
    \[|\psi\rangle = \dfrac{1}{\sqrt{2}}|0\rangle|0\rangle^{\otimes q} + \dfrac{1}{\sqrt{2}}|1\rangle|0\rangle^{\otimes q}\]
    Where the first register corresponds to the ancilla qubit and the second register corresponds to the data qubits.

    A CNOT gate controlled by the ancilla qubit acts on the first active input qubit (index $q - n + 1$) of the data register, creating entanglement and the resultant state:
    \[|\psi\rangle = \dfrac{1}{\sqrt{2}}|0\rangle|0\rangle^{\otimes q} + \dfrac{1}{\sqrt{2}}|1\rangle|\mathbf{e}_{q-n+1}\rangle\]
    where $|\mathbf{e}_{q-n+1}\rangle$ is a unary vector of size $q$ with the $(q - n + 1)$-th qubit set to $|1\rangle$.
    
    \item \textbf{Amplitude encoding in a unary subspace} — Given an input vector $\mathbf{x}$ and the qubits initialized in the unary state $|\mathbf{e}_{q-n+1}\rangle$, we encode $\mathbf{x}$ in the Hilbert space as:
    \[S(\mathbf{x})|\mathbf{e}_{q-n+1}\rangle = x_1 |\mathbf{e}_{q-n+1}\rangle + x_2 |\mathbf{e}_{q-n+2}\rangle + \dots + x_n |\mathbf{e}_{q}\rangle = \sum_{i=1}^n x_i|\mathbf{e}_{q-n+i}\rangle\]
    where $S(\mathbf{x})$ is a data-dependent encoding unitary operating on the bottom $n$ qubits.

    Because the unitary $S(\mathbf{x})$ is constructed entirely from Reconfigurable Beam Splitter (RBS) gates, which preserve the $|00\rangle$ state, it has no effect on the ground state ($S(\mathbf{x})|0\rangle^{\otimes q} = |0\rangle^{\otimes q}$). Consequently, applying the operator $I \otimes S(\mathbf{x})$ to the full system transforms the overall state to:
    \[|\psi\rangle = \dfrac{1}{\sqrt{2}}|0\rangle|0\rangle^{\otimes q} + \dfrac{1}{\sqrt{2}}|1\rangle\sum_{i=1}^{n} x_i|\mathbf{e}_{q-n+i}\rangle\]

    \item \textbf{Pyramidal orthogonal transformation} — Any orthogonal matrix $\mathbf{W}$ in $\mathbf{y} = \mathbf{Wx}$ can be decomposed into a network of RBS gates under nearest-neighbour connectivity. This yields a parameterized unitary $U(\boldsymbol{\theta})$ that maps the $n$-dimensional input subspace to the $m$-dimensional output subspace:
    \[U(\boldsymbol{\theta})S(\mathbf{x})|\mathbf{e}_{q-n+1}\rangle = \sum_{j=1}^m \left( \sum_{i=1}^n {W}_{ji}x_{i} \right) |\mathbf{e}_{q-m+j}\rangle\]
    Because the parameterized unitary $U(\boldsymbol{\theta})$ is also constructed entirely from RBS gates, it preserves the ground state ($U(\boldsymbol{\theta})|0\rangle^{\otimes q} = |0\rangle^{\otimes q}$). Applying the operator $I \otimes U(\boldsymbol{\theta})$ to the joint system yields the overall state:\[
    |\psi\rangle = \dfrac{1}{\sqrt{2}}|0\rangle|0\rangle^{\otimes q} + \dfrac{1}{\sqrt{2}}|1\rangle\sum_{j=1}^m \left( \sum_{i=1}^n {W}_{ji}x_{i} \right)|\mathbf{e}_{q-m+j}\rangle\]
    
    \item \textbf{Tomography and sign extraction} — Finally, we must recover the transformed components $\sum_{i=1}^n W_{ji}x_i$ from the circuit. To do so, we introduce $S(\mathbf{r})$, an encoding unitary parameterized to load the uniform norm-1 vector $\mathbf{r} = (\frac{1}{\sqrt{q}}, \frac{1}{\sqrt{q}}, \dots , \frac{1}{\sqrt{q}})$ onto the full $q$-dimensional subspace. First, we apply the adjoint operator $I \otimes S^\dagger(\mathbf{r})$ to the full system.
    \[|\psi\rangle = \dfrac{1}{\sqrt{2}}|0\rangle|0\rangle^{\otimes q} + \dfrac{1}{\sqrt{2}}|1\rangle S^\dagger(\mathbf{r})\sum_{j=1}^{m}\left( \sum_{i=1}^n {W}_{ji}x_{i} \right)|\mathbf{e}_{q-m+j}\rangle\]
    Second, we flip the ancilla qubit with an X gate and perform a CNOT gate on the first active qubit (index $1$) of the data register, controlled by the ancilla qubit.
    \[|\psi\rangle = \dfrac{1}{\sqrt{2}}|1\rangle|\mathbf{e}_{1}\rangle + \dfrac{1}{\sqrt{2}}|0\rangle S^\dagger(\mathbf{r})\sum_{j=1}^{m} \left( \sum_{i=1}^n {W}_{ji}x_{i} \right) |\mathbf{e}_{q-m+j}\rangle\]
    Third, we apply $I \otimes S(\mathbf{r})$ to the full system. On the $|1\rangle$ branch, this loads the uniform vector $\mathbf{r}$. On the $|0\rangle$ branch, $S(\mathbf{r})$ and $S^\dagger(\mathbf{r})$ perfectly cancel out ($S(\mathbf{r})S^\dagger(\mathbf{r}) = I$), restoring the transformed data components.
    \[|\psi\rangle = \dfrac{1}{\sqrt{2}}|1\rangle\sum_{k=1}^q \dfrac{1}{\sqrt{q}}|\mathbf{e}_{k}\rangle + \dfrac{1}{\sqrt{2}}|0\rangle\sum_{j=1}^{m} \left( \sum_{i=1}^n {W}_{ji}x_{i} \right) |\mathbf{e}_{q-m+j}\rangle\]
    Finally, a Hadamard gate on the ancilla qubit mixes the states to provide the result:
    \[|\psi\rangle = \dfrac{1}{2}|0\rangle \left( \sum_{j=1}^m y_j|\mathbf{e}_{q-m+j}\rangle + \sum_{k=1}^q \dfrac{1}{\sqrt{q}}|\mathbf{e}_k\rangle \right) + \dfrac{1}{2}|1\rangle \left( \sum_{j=1}^m y_j|\mathbf{e}_{q-m+j}\rangle - \sum_{k=1}^q \dfrac{1}{\sqrt{q}}|\mathbf{e}_k\rangle \right)\]
    where $y_j = \sum_{i=1}^n W_{ji}x_i$. By isolating the measurements on the specific $m$-dimensional subspace, the $j$-th component of the output vector can then be inferred using the formula:
    \begin{align*}
        \sqrt{q}(\Pr[0, \mathbf{e}_{q-m+j}] - \Pr[1, \mathbf{e}_{q-m+j}]) &= \dfrac{\sqrt{q}}{4}\left(y_j + \dfrac{1}{\sqrt{q}}\right)^2 - \dfrac{\sqrt{q}}{4}\left(y_j - \dfrac{1}{\sqrt{q}}\right)^2
        = y_j
    \end{align*}
    
\end{enumerate}

\clearpage
\section{Nomenclature}
\label{sec:nomenclature}
\subsection*{Data and Operator Formulation}
\begin{description}
    \item[$\mathcal{G}$] Target operator mapping input functions to pointwise predictions.
    \item[$\mathcal{D}$] Supervised training dataset.
    \item[$\mathbf{u}_i, \mathbf{u}$] Discretized input function evaluated at fixed sensor locations.
    \item[$d_u$] Dimensionality of the discretized input function.
    \item[$\mathbf{y}_{ij}, \mathbf{y}$] Spatial or temporal query coordinate.
    \item[$d_y$] Dimensionality of the query coordinate.
    \item[$s_{ij}$] Ground-truth value of the operator evaluated at a specific coordinate.
    \item[$N$] Total number of input trajectories or hypothetical scenarios in a batch.
    \item[$M, M_i$] Maximum and specific number of spatial/temporal query coordinates evaluated per scenario.
\end{description}

\subsection*{Network Architecture}
\begin{description}
    \item[$\mathbf{b}$] $p$-dimensional branch network output vector.
    \item[$\mathbf{t}$] $p$-dimensional trunk network output vector.
    \item[$\theta$] Trainable parameters of the sub-networks, representing Reconfigurable Beam Splitter rotation angles.
    \item[$d_\theta$] Total number of trainable parameters.
    \item[$n$] Hidden layer width and input vector dimension for orthogonal linear transformations.
    \item[$\mathbf{W}$] Orthogonal weight matrix emulated by the parameterized quantum circuit.
\end{description}

\subsection*{Quantum Circuit}
\begin{description}
    \item[$L$] Number of independent quantum DeepONet models in the ensemble.
    \item[$S(\cdot)$] Data loader parameterized unitary circuit utilizing amplitude encoding.
    \item[$U(\cdot)$] Pyramidal parameterized unitary circuit executing orthogonal linear transformations.
    \item[$|\mathbf{e}_i\rangle$] Unary computational basis state within the quantum register.
    \item[$N_{shots}$] Total number of quantum measurement shots budgeted per inference.
    \item[$\lambda$] Noise strength parameter for simulated depolarizing channels.
\end{description}

\subsection*{Uncertainty Quantification}
\begin{description}
    \item[$\mu(\mathbf{u})(\mathbf{y})$] Mean prediction output calculated across the ensemble.
    \item[$\sigma(\mathbf{u})(\mathbf{y})$] Ensemble standard deviation quantifying epistemic disagreement.
    \item[$r_{ij}$] Adaptive nonconformity score evaluating prediction error normalized by local variance.
    \item[$\hat{q}$] Empirical quantile of the nonconformity scores derived from calibration data.
    \item[$C_\alpha(\mathbf{u}, \mathbf{y})$] Adaptive, distribution-free prediction set.
    \item[$\alpha$] User-specified miscoverage rate for the conformal prediction intervals.
\end{description}


\newpage
\section{Training hyperparameters for Ideal Simulations}\label{sec:train_hp}

We list the hyperparameters used during training in Table \ref{tab:training_hyperparams}. We utilized Lambda decay to reduce the learning rate.

\begin{table}[t]

\caption{Training hyperparameters. ``Full'' indicates full-batch gradient descent. $\gamma$ is the decay factor. ``--'' indicates a constant learning rate (no scheduler). Power-system rows report train/cal/test split fractions}
\label{tab:training_hyperparams}

\begin{center}
\small
\begin{tabular}{lccccccccc}
\hline
\textbf{Experiment} & $N_{train}$ & $N_{cal}$ & $N_{test}$ & \textbf{Batch} & \textbf{Iters} & \textbf{LR} & \textbf{Min LR} & $\gamma$ & \textbf{Loss} \\
\hline

Antiderivative & 200  & 50  & 50  & Full & 30k & $1\mathrm{e}{-3}$ & --              & --   & MSE \\
Advection      & 1000 & 200 & 200 & Full & 40k & $1\mathrm{e}{-3}$ & $5\mathrm{e}{-4}$ & 0.99 & MSE \\
Offline V-to-V & 0.8  & 0.1 & 0.1 & 256  & 40k & $5\mathrm{e}{-4}$ & --              & --   & MSE \\
Offline V-to-P & 0.8  & 0.1 & 0.1 & 64   & 40k & $5\mathrm{e}{-3}$ & $5\mathrm{e}{-4}$ & 0.99 & Rel. $L_2$ \\
Online V-to-V  & 0.8  & 0.1 & 0.1 & 256  & 15k & $1\mathrm{e}{-2}$ & --              & --   & MSE \\

\hline
\end{tabular}
\end{center}

\end{table}

\section{Multinomial sampling vs per-shot simulations}
\label{sec:multinomial}


To accelerate the simulation of quantum circuits, we employ multinomial sampling from the final statevector or density matrix probability distribution, rather than simulating each shot individually. This single simulation followed by sampling is computationally efficient.

To validate this approach, we compare its results against a full per-shot simulation, which applies noise at every step for every shot. Figure \ref{fig:appendix_c_compare} shows this comparison for the antiderivative operator. The results for all metrics (L2 Error, Coverage, Average Width, and Max Width) demonstrate correlation, validating that our multinomial sampling method is an accurate and efficient substitute for the more costly per-shot simulation.

\begin{figure}[t]
\centering
\includegraphics[width=1.0\textwidth]{figures/appendix_c_parity_plot.pdf}
\vspace{-2em}
\caption{Validation of multinomial sampling. Parity plots comparing the outputs of per-shot simulation (x-axis) against multinomial sampling (y-axis). The dashed line represents $y=x$.}
\label{fig:appendix_c_compare}
\end{figure}



\clearpage
\section{System snapshots}\label{sec:snapshots}

The realistic noise model was constructed based on calibration data from three Qiskit backends: \texttt{fake\_brisbane}, \texttt{fake\_torino}, and \texttt{fake\_marrakesh}. The noise model was generated using these properties. All simulations were transpiled to the basis gate set and coupling map of the respective backend. 

\begin{table}[htbp]

\caption{Backend characteristics for the simulated quantum devices. Reported values include hardware configuration, coherence times ($T_1$, $T_2$), readout errors, and gate errors. All values are averaged, minimum, and maximum across qubits unless otherwise specified.}
\label{tab:backend_specs}

\begin{center}
\small

% --- Table 1 ---
\begin{tabular}{lccc}
\hline
\textbf{Backend} & \textbf{Qubits} & \textbf{Basis Gates} & \textbf{Coupling} \\
\hline

fake\_brisbane  & 127 & \{ecr, id, rz, sx, x\} & 144 edges \\
fake\_torino    & 133 & \{cz, id, rz, sx, x\}  & 300 edges \\
fake\_marrakesh & 156 & \{cz, id, rz, sx, x\}  & 352 edges \\

\hline
\end{tabular}

\vspace{0.6em}

% --- Table 2 ---
\begin{tabular}{lcccccc}
\hline
\textbf{Backend} & $T_1$ (avg) & $T_1$ (min) & $T_1$ (max) & $T_2$ (avg) & $T_2$ (min) & $T_2$ (max) \\
\hline

fake\_brisbane  & $2.34\mathrm{e}8$ & $9.94\mathrm{e}6$ & $4.25\mathrm{e}8$ & $1.61\mathrm{e}8$ & $1.71\mathrm{e}7$ & $4.12\mathrm{e}8$ \\
fake\_torino    & $1.74\mathrm{e}8$ & $3.16\mathrm{e}6$ & $3.08\mathrm{e}8$ & $1.45\mathrm{e}8$ & $6.13\mathrm{e}6$ & $3.50\mathrm{e}8$ \\
fake\_marrakesh & $2.07\mathrm{e}8$ & $6.80\mathrm{e}6$ & $4.99\mathrm{e}8$ & $1.47\mathrm{e}8$ & $1.80\mathrm{e}5$ & $5.35\mathrm{e}8$ \\

\hline
\end{tabular}

\vspace{0.6em}

% --- Table 3 ---
\begin{tabular}{lccc}
\hline
\textbf{Backend} & \textbf{Readout Err (avg)} & \textbf{Readout Err (min)} & \textbf{Readout Err (max)} \\
\hline

fake\_brisbane  & 3.11 & 0.56 & 23.95 \\
fake\_torino    & 4.67 & 0.51 & 56.52 \\
fake\_marrakesh & 2.70 & 0.27 & 47.39 \\

\hline
\end{tabular}

\vspace{0.6em}

% --- Table 4 ---
\begin{tabular}{lcccc}
\hline
\textbf{Backend} & \textbf{1Q Err (avg)} & \textbf{1Q Err (min)} & \textbf{1Q Err (max)} & \textbf{2Q Err (avg)} \\
\hline

fake\_brisbane  & 0.11 & 0.00 & 14.09 & 1.96 \\
fake\_torino    & 0.06 & 0.00 & 1.74  & 8.00 \\
fake\_marrakesh & 0.06 & 0.00 & 4.20  & 8.06 \\

\hline
\end{tabular}

\end{center}

\end{table}



\end{document}


